\setchapterpreamble[u]{\margintoc}
\chapter{A Map of Learning Problems}
\labch{map-of-learning-problems}

\sources{Zhang, Lipton, Li and Smola, \emph{Dive into Deep Learning}
\cite{d2l2023}, Chapter~1.}

With the four components of \refch{learning-problem} in place, the standard
catalogue of problem types can be laid out as variations on them: what the data
contain, what is being predicted, and whether the learner influences what it
sees next.

\section{Supervised learning}
\labsec{supervised}

Every example carries a label, and the goal is to model the conditional
distribution of the label given the features --- the object \(f_{Y\mid X}\) of
\refsec{joint}.

\subsection{Regression}
\labsubsec{regression}

The label is a real number and the natural loss is squared error, so by
\refsec{conditional-expectation} the fitted function approximates
\(\E[Y\mid X]\). Everything in \refch{orthogonality} applies verbatim when the
model is linear.

\subsection{Classification}
\labsubsec{classification}

The label is one of finitely many \emph{categories} or \emph{classes}. Two
classes make it binary classification; more make it multiclass. The output is a
probability vector, and the loss is cross-entropy, which
\refsec{losses-from-likelihood} derives rather than assumes.

\emph{Hierarchical} classification exploits structure among the classes so that
confusing two neighbouring categories is penalised less than confusing two
distant ones. \emph{Multi-label} classification drops the assumption that the
categories are mutually exclusive, and models each label with its own Bernoulli
output.

\marginnote{Accuracy is a poor summary when classes are unbalanced, for the
reason computed in \refsec{bayes}: with a rare positive class, most positive
predictions can be wrong even at very high accuracy.}

\subsection{Search and ranking}
\labsubsec{ranking}

The output is an ordering of a candidate set rather than a single value, and
quality depends on which items appear near the top. Losses here are built to
approximate rank-based measures, which are themselves not differentiable --- an
instance of the surrogate problem noted in \refsubsec{objective}.

\subsection{Recommender systems}
\labsubsec{recommenders}

Predict a user's response to an item from a sparse, partially observed matrix of
past interactions. The classical treatment is a low-rank factorisation of that
matrix, which is \refsec{low-rank} applied to a table of ratings.

\begin{remark}
Recommendation is the clearest case in which the fitted system changes the data
it will later be trained on: what is recommended is what gets seen, and what
gets seen is what gets recorded. Censoring, incentives, and feedback loops are
open research problems, not solved engineering details.
\end{remark}

\subsection{Sequence learning}
\labsubsec{sequences}

Input and output are sequences whose lengths need not match. Tagging and parsing
attach structure to each position; automatic speech recognition maps audio to
text; text-to-speech maps text to audio; machine translation maps a sequence in
one language to another, generally of different length and different word order.

\section{Unsupervised learning}
\labsec{unsupervised}

No labels are supplied, and the goal is to describe the structure of the data
itself.

\begin{description}
	\item[Clustering] group examples so that members of a group resemble each
	other more than they resemble members of other groups.
	\item[Principal component analysis] find the few directions along which the
	data vary most --- exactly \refsec{pca}.
	\item[Causality and probabilistic graphical models] describe which variables
	influence which, rather than merely which co-vary.
	\item[Generative adversarial networks] fit a mechanism that produces samples
	resembling the data, by training a generator against a discriminator.
\end{description}

\marginnote{Only the second of these is fully settled by linear algebra.
Correlation is what a covariance matrix records, and \refsec{covariance-def}
notes that it detects linear association only --- causal structure is not
recoverable from it without further assumptions.}

\section{Interacting with an environment}
\labsec{environment}

The problems above are \emph{offline}: a fixed dataset is collected first, the
model is fitted second, and the model never influences what it is shown. This is
convenient and often unrealistic.

\begin{definition}
\emph{Distribution shift} occurs when the distribution generating the data at
deployment differs from the one that generated the training data.
\end{definition}

Since \refsec{lln} is what licensed the transfer from sample to population, and
it assumes a single fixed distribution, shift is not a small perturbation of the
theory --- it removes the hypothesis the theory rests on.

\section{Reinforcement learning}
\labsec{reinforcement}

Here the learner is an \emph{agent} placed in a loop with an environment. At
each step it receives an \emph{observation}, chooses an \emph{action} according
to its \emph{policy}, and receives a \emph{reward}; the action influences what is
observed next.

Two difficulties distinguish this from supervised learning. The reward may
arrive long after the action responsible for it, so the agent faces the
\emph{credit assignment} problem. And the agent must balance exploiting what it
already believes to be good against exploring alternatives it has not tried.

\begin{description}
	\item[Markov decision process] the standard model, in which the next state
	depends only on the present state and the chosen action.
	\item[Multi-armed bandit] the one-state case, in which the agent chooses
	repeatedly among fixed options with unknown rewards.
	\item[Contextual bandit] a bandit in which an observation is supplied before
	each choice, but actions still do not affect future states.
\end{description}

\marginnote{The bandit cases are where the estimation machinery of
\refch{estimation} applies most directly: each arm's value is a mean to be
estimated, and the exploration problem is about how confident those estimates
are.}

\section{Terminology carried forward}
\labsec{terminology}

A last group of words appears throughout the literature and is worth fixing
before Part~V: \emph{estimation}; \emph{neural network}; \emph{layer};
\emph{kernel method}; \emph{decision tree}; \emph{graphical model}; and
\emph{representation learning} --- the idea that the useful features should be
learned along with the predictor rather than designed in advance. That last term
is the one \refch{column-space-to-layers} takes up.
